\section{Social claims after source and recipient checks}\label{sec:conduct}

The corpus contains 6,882 nonempty public-message occurrences alongside 526 proposals and 113 settlements. The selected incidents include an explicit omission rationale, repeated assurances withdrawn under financial distress, an appeal to withhold a rival's cash, a false obstacle contradicted by an applied private note, and a public bargaining limit below the privately stated limit. We present the exchanges with their context and outcomes, rather than infer behavior from isolated quotations. These are case studies, not a deception rate or a moral ranking of the participants.

% Preserved normative and evidence protocol from the pre-transcript draft.
\subsection{The normative contract and the evidence chain}\label{sec:conduct-norm}

The standard system prompt tells each agent to ``win the game by maximizing long-term advantage and being the last non-bankrupt player'' and permits ``any legal strategy or personality (aggressive, deceptive, cooperative, conservative, manipulative, opportunistic, friendly, hostile, etc.)''. It describes the public message as a place to ``negotiate, bluff, deceive, cooperate, intimidate, joke, or stay silent'', and states: ``Deception in public chat is allowed; private thoughts should reflect true reasoning.'' The action-only retry, used only after a provider-refused first attempt, removes the private-note field and shortens that sentence to ``Deception in public chat is allowed.'' It also instructs: ``Do not include private\_thought or a reasoning transcript.'' This is a response-format exclusion, not a separate honesty requirement for provider-returned explanations; the records do not establish which provider fields that instruction controls. The win and legal-action instructions appear in every prompt. Public deception is permitted throughout, while an applied private note is asked to reflect true reasoning only under the standard condition. Every Opus first attempt from game~5 onward was refused, so its applied decisions there were resolved under the action-only condition or by fallback (Section~\ref{sec:setup}). Opus therefore mostly played without the private-honesty clause, not merely without a private note.

The contract fixes how the rest of this section is labelled. A false public statement is not a violation of these instructions, and by itself it is not evidence that an agent is misaligned with them. A false private note written under the standard condition is a candidate breach of a request, but error, stale context and disagreement between channels remain live explanations, so a false note does not on its own establish dishonesty. Rule-breaking actions cannot be applied, because the engine executes only menu actions; invalid attempts are recorded instead. We describe the deceptive and adversarial tactics the records support while distinguishing them from violations of the assigned task. The analysis asks how claims were produced, how recipients detected and answered them, and what changed. This differs from AgentMonopoly, whose prompt encourages bluffing but prohibits public messages that contradict the submitted action \citep{wu2026agentmonopoly}, and from the Vending-Bench Arena communication study, which applies an external taxonomy to agents given no conduct instruction \citep{li2026arena}. An observed contradiction can be a measured violation in one of those environments and an authorized move in ours.

We distinguish four objects. A \emph{factual assertion} concerns the board or rules. An \emph{option omission} leaves a relevant legal route unstated, potentially creating a false necessity implication. A \emph{bargaining posture}, such as ``final'', expresses a stated limit rather than an observed reservation price. An \emph{assurance} concerns future behavior, whose obligation and expiry may be unspecified. State and engine checks can settle some factual or feasibility propositions; they cannot alone decide intent or whether an omission changes beliefs. Speaker-observation checks distinguish an error from a statement contradicted by information available to its author \citep{turpin2023unfaithful,lanham2023faithfulness}.

For each case we link the proposition, speaker channels, recipient's actual prompt, response, subsequent events and legal alternatives. Exposure is checked at the response where reliance is proposed. Information defects remain part of that condition: G4 d148--152 contains reversed acceptances, misattributed rejections and self-partner counters; G13 d252--253 contains a self-partner counter. The inspected G6 d19--21, G11 d255--258 and G19 d135--136 timelines contain none of these known formatter manifestations. These are exact-row checks, not assertions that other prompts are error-free or that a defect caused a response.



\subsection{Five conduct chronologies}
\label{sec:conduct-framing}\label{sec:conduct-assurance}\label{sec:conduct-thirdparty}

The cards below are verbatim excerpts, with separate labels for public messages, applied private notes and returned provider explanations. Surrounding prose supplies the board context, actual action, recipient exposure and later consequence. A bracketed ellipsis marks omitted text; errors within quotations are left intact. The presentation follows the useful context--excerpt--interpretation structure of published agent case studies \citep{li2026arena}, while preserving the different channels and contract in this benchmark. Appendix~\ref{app:transcript-provenance} describes selection and source checks.

% Contextual transcript exhibits. Quote bodies are generated separately from
% exact original-field spans; all prose here is the authors' interpretation.

\Needspace{18\baselineskip}
\TranscriptTitle{omission}{An insurance pitch that omits available collateral}

Grok wants cash without giving up its yellow properties. Opus wants to prevent a rival from inheriting them. In G13, Grok has \$35 and offers its mortgaged Pennsylvania Avenue to Opus for \$200. Sol owns all four railroads, including Short Line, whose \$200 rent threatens Grok. Ventnor, one of Grok's yellows, is still clear and can be mortgaged for \$130. Grok first presents the sale as insurance against an estate transfer:

\TranscriptQuote{g13-opening-insurance}

Opus receives this appeal and counters at \$130 in d251. Its returned explanation values keeping Grok solvent but calls the \$200 price too high. That rationale follows the seller's opening appeal; it is not evidence of an independently formed motive. Grok answers:

\TranscriptQuote{g13-seller-counter}

The cash-only shortfall is real: \$35 plus \$130 is \$35 below the rent. The missing fact is Grok's own collateral. The same response records both the mortgage route and a preference to avoid it:

\TranscriptQuote{g13-seller-floor}

Its returned provider explanation is more explicit about the presentation:

\TranscriptQuote{g13-omission-explanation}

This is an excerpt from a 21,134-character explanation that also contains conflicting arithmetic and ownership statements. It is a recorded rationale for concealment, not a warrant to treat the entire explanation as one stable intention. Opus sees the new pitch and accepts:

\TranscriptQuote{g13-acceptance}

\TranscriptFinding{Events 1808--1812 transfer Pennsylvania to Opus, \$200 to Grok and \$16 of buyer mortgage interest to the bank. Grok reaches \$235. In an exact-prefix alternative, accepting \$130 returns control to Grok's management menu; mortgaging Ventnor then reaches \$295, enough for the \$200 rent with \$95 left. The eventual outcome is also relevant: at d311, turn 150, Grok owes Sol \$200 on Short Line, has \$49 and no financing action, and loses the yellows to Sol in bankruptcy. The earlier sale did not prevent that later estate transfer.}

\TranscriptReading{This is the strongest retained example of a delivered omission with an explicit concealment rationale. The financial-necessity reading leaves out a feasible way to pay. A narrower refusal-posture reading also remains: Grok may be threatening not to accept \$130, rather than asserting that accepting it makes payment impossible. Its note sets a \$165 floor to preserve Ventnor's rent. The returned provider explanation accompanying Opus's acceptance repeats the estate-denial stakes but never says the cheaper route is impossible. The \$70 price difference is observed; a \$70 persuasion effect is not identified. The opening 5/36 probability is itself wrong---Short Line is 4/36 from Atlantic. An earlier claim that inheritance would complete Opus's greens was publicly corrected at d244 because North Carolina remained unowned; that correction was also in Opus's prompt.}

\TranscriptSource{G13 d250--253, turn 107; follow-up d311, turn 150. Grok's quoted response is standard. Opus's d251 and d253 selected responses use action-only recovery and have no applied private note. The d252--253 timelines include a self-partner counter description; correct holdings remain available. Exact spans and source hashes are retained in the transcript manifest (Appendix~\ref{app:transcript-provenance}).}

\TranscriptTitle{withdrawal}{A standing offer withdrawn when its recipient is nearly broke}

In G18, Opus wants Sol's mortgaged Pacific Avenue to complete its greens. Sol treats Pacific as a valuable blocker. Opus offers \$250 at d421, calls it final, and repeatedly says the offer stands at d423, d425 and d427. After collecting \$750 rent from Sol, it repeats the assurance:

\TranscriptQuote{g18-assurance}

Sol initially asks for more. Then another \$750 rent payment forces it to sell orange hotels and houses and mortgage its remaining reds at d432--436. It has \$49 afterward. Its subsequent \$450 and \$350 packages do not settle. At d441, Sol invokes the advertised price for Pacific alone:

\TranscriptQuote{g18-invocation}

Opus rejects it:

\TranscriptQuote{g18-withdrawal}

The selected response's returned explanation gives a direct reason to withhold the money:

\TranscriptQuote{g18-withdrawal-explanation}

\TranscriptFinding{The d431 assurance is present in both players' delivered memories at invocation and withdrawal. No sale occurs at \$250. Later Opus offers \$150 and then \$225; Sol continues to ask more, without agreement. Sol's invoked price therefore becomes unavailable precisely when its finances have deteriorated.}

\TranscriptReading{The record shows repeated assurance, explicit invocation and opportunistic withdrawal, with the buyer citing the seller's near-insolvency. It does not show that Sol incurred a cost in reliance on the assurance: the intervening liquidation pays a second rent bill, and its earlier proposals sought more than \$250. Nor does withdrawal prove dishonesty when the offer was made. The game has no standing-offer contract object or specified expiry, and the instruction to remain consistent allows a reason to change. Those limits distinguish an observed bargaining reversal from a demonstrated broken binding commitment \citep{empty_commitments_v2,commitment_standing}.}

\TranscriptSource{G18 d421--442, especially d431, d441 and d442. Opus's explanation is provider-returned text under action-only recovery, not an applied private note. No costly-reliance or winner counterfactual is claimed.}

\TranscriptTitle{boycott}{An appeal to withhold a rival's building cash}

In G20, Opus has sold Illinois to Grok for \$500, completing Grok's reds. Gemini sees a complete rival set with little cash to develop it. At d96 it asks the other players to stop financing Grok through purchases:

\TranscriptQuote{g20-coordination-appeal}

Grok's next applied private note recognizes the appeal:

\TranscriptQuote{g20-recognition-note}

That d97 received action is \End, despite text describing financing and building. Three red houses are actually built at d98; Grok mortgages New York at d99 and then offers Vermont to Gemini for \$150:

\TranscriptQuote{g20-sale-response}

Gemini refuses:

\TranscriptQuote{g20-concrete-refusal}

\TranscriptFinding{Gemini's appeal is followed by its own rejection of a concrete cash purchase. Grok instead mortgages Vermont at d102. The exchange is an observable attempt to coordinate financial denial, and the target explicitly calls it a boycott. It is not merely an analyst's description of parallel play.}

\TranscriptReading{The evidence does not establish an agreed cartel. Sol had already stated its own reason not to finance Grok at d95, before Gemini's appeal. Opus later proposes cash purchases from Grok twice, neither of which settles. Thus the initiator follows its stated position, but the record does not show the table committing to it or identify the appeal's effect on another player's choice. Grok's ``not a house fund'' framing promotes its own offer; it does not earmark how the cash can be spent.}

\TranscriptSource{G20 d95--102. Gemini's d96 appeal is delivered to Grok at d97; its d101 refusal answers the actual d100 proposal. The quoted d97 note accompanies selected attempt 1 after a malformed attempt 0; its descriptions are not applied construction.}

\TranscriptTitle{stock}{A building obstacle contradicted by the speaker's own note}

Late in G3, Gemini has a dark-blue group and Grok is discussing the building threat. Only one bank house remains. Grok publicly says Gemini cannot add the Park Place house unless somebody frees buildings, but its applied note identifies that one existing house as enough and lists mortgage routes:

\TranscriptQuote{G3-2}

The same response's private note says:

\TranscriptQuote{G3-3}

Gemini receives the public claim. Its next applied note gives a reason to hold:

\TranscriptQuote{G3-5}

\TranscriptFinding{Gemini ends its next management decision without building. At that actual d364 opportunity, after an intervening \$15 tax, it has \$43. An exact-prefix legal sequence mortgages Kentucky for \$110 and Atlantic for \$130, then buys the \$200 Park Place house, leaving \$83 and exhausting the one remaining bank house. No other player has to release a building.}

\TranscriptReading{The asserted need for additional stock is false for this local route, while the cash-only shortfall is true. The speaker's applied note records a contrary route, but its longer provider explanation repeatedly treats a release as necessary, so the channels do not establish one coherent deliberate deception. Gemini's earlier d362 hold already precedes the claim; its d364 note favors retaining unmortgaged income. Exposure followed by another hold therefore does not identify reliance on the false obstacle. This is misleading pressure with conflicting explanatory evidence, not a demonstrated induced error.}

\TranscriptSource{G3 d362--364. The retained exact-prefix legality audit supplies the financing witness: mortgage Kentucky and Atlantic, then buy one Park Place house. The bounded transcript review independently reproduced that route. The endpoint is this legal purchase, not a prediction of survival or the winner.}

\TranscriptTitle{final}{A public final offer below the private limit}

Grok wants North Carolina Avenue from Sol to advance a green set. It publicly presents \$600 as the end of the negotiation:

\TranscriptQuote{g7_d42_public}

The same response's private field leaves room to pay more:

\TranscriptQuote{g7_d42_private}

Sol accepts the public figure:

\TranscriptQuote{g7_d43_public}

\TranscriptFinding{Events 324--327 transfer North Carolina from Sol to Grok and \$600 in the other direction. Sol's applied note says it would rather reinvest in cheaper sets than hold out for another \$25. The private \$625 contingency is never tested.}

\TranscriptReading{This is a direct public/private discrepancy about a bargaining limit, followed by settlement at the advertised limit. It is a clearer bluff than a false board statement shared across all channels, but it is a lower-stakes incident than the solvency pitch or withdrawal above. Nothing in the recipient's stated reason establishes that ``final'' secured the price. Bluffing was expressly permitted by the benchmark prompt.}

\TranscriptSource{G7 d42--43, turn 25; both selected attempts are standard and valid at index 0. Sol's actual d43 prompt contains the \$600 final-offer message.}


\subsection{Counterexamples and unresolved uptake}

The selected alarming statement is not always the recipient's belief. In G19, Sol explicitly corrects a false color claim and accepts the transaction anyway (Transcript~\ref{tr:abort}). In G6 d19--21, Sol's applied note similarly rejects Grok's false claim about an orange monopoly while accepting St.~Charles, Pennsylvania Railroad and \$400 for Vermont. The proposer later calls that trade an overpay. Both G6 and G19 also contain content/tool conflicts about whether to propose at all. Acceptance does not by itself establish either uptake of the premise or a coherent proposer intention.

G4 supplies a different contrast. Grok's false \$200 yellow-house cost is followed by Gemini's public ``Fair point'' and a changed request. But Gemini's d148 note already records the true \$150 figure and the contingency of taking yellow with compensation; its d150 returned explanation retains \$150. The final exchange completes groups for both parties at a \$50 payment, after a \$143 request. This does not identify a belief change or a \$93 concession caused by the false figure. In G11 d256, the recipient's explanation transiently echoes a false completion claim, then checks the missing St.~Charles and corrects itself before rejecting. The subsequent true completion offer is also rejected. Appendix~\ref{app:social-detail} gives the surrounding chronologies.

Table~\ref{tab:framing-cases} summarizes five previously adjudicated proposition cases. G13 leaves belief in the omitted-route necessity unresolved; G6 and G19 contain explicit correction with acceptance; G4 has prior correct knowledge; G11 has echo followed by correction and rejection. These observations refute particular sufficiency rules, not every influence hypothesis. Controlled negotiation studies can measure intervention effects \citep{bianchi2024negotiationarena}; this cohort has neither randomized message variants nor repeated recipient continuations.

\begin{table}[t]
\centering\small
\caption{Five original-checked chronologies of false or misleading propositions. ``Recipient record'' refers to the recipient's channels at or before its response. Selected contrasts, not prevalence.}
\label{tab:framing-cases}
\begin{tabularx}{\linewidth}{@{}L{1.5cm}YYYL{2.0cm}@{}}
\toprule
Case & Proposition and its status & Speaker records & Recipient record and response & Outcome\\
\midrule
\decision{13}{252} & ``\$130 leaves me \$35 short''; true in cash, omits a \$130 mortgage (\$295 reachable) & Applied note acknowledges the route; returned reasoning states the omission and also contradicts itself & Solvency value noted at d251 after seller's d250 appeal; silent on the \$130 route; accepts \$200 & Settled at \$200\\
\decision{19}{135} & ``States is the third orange''; false & Applied note says abort; content and reasoning select \End; tool proposes & Corrects the premise publicly; note treats States as pink; accepts & Settled\\
\decision{11}{255} & ``Virginia completes your pinks''; false & Applied note repeats the premise; later concedes the mistake & Reasoning echoes, then corrects; applied note rejects & Rejected twice\\
\decision{4}{149} & ``\$200 a house'' for yellow; false (\$150) & Applied note states a one-house-each limit without the \$200 figure & True \$150 and a contingency recorded at d148; public ``Fair point''; switches & Settled at \$50\\
\decision{6}{19} & ``Your oranges are the board's only set''; false & Applied note repeats the premise; content and reasoning select \End; tool proposes & Applied note rejects the premise; accepts & Settled; speaker's later note calls it an overpay\\
\bottomrule
\end{tabularx}
\end{table}


\subsection{Disclosure and conflicting channels}\label{sec:conduct-channels}

Three further transcripts preserve consequential events that should not be relabelled as malicious plans. G19's sent proposal carries a note saying to abort it; G12's received bankruptcy call accompanies content requesting a sale and explanations arguing for survival; G20's public field discloses a private note even though assistant content requested an empty message (Appendix~\ref{app:transcript-controls}). The first two change assets or survival. The third is delivered and recognized by an opponent. None supports an unqualified claim of intentional betrayal, surrender or disclosure.

Grok was the only configuration run with required tool choice. It has 139 action-name conflicts between assistant content and the received call. In fifteen, the parsed argument objects are equal and only the function name differs; five of those fifteen also have literally equal JSON substrings. Four of the fifteen accept trades while their messages argue for rejection. All 139 selected conflicts already occur in saved nonstreamed response bodies, excluding local stream assembly for these checked records. Generation, provider processing and OpenRouter remain unresolved sources (Section~\ref{sec:validation}). These are action-attribution findings, not a deception measure.

A mechanical scan for overlaps of at least forty characters between each applied private note and its public message finds seven overlaps across all 6,989 decisions, fewer of which carry both fields. Six restate public-state content; one is the G20 structural disclosure (Transcript~\ref{tr:disclosure}). Paraphrased disclosure would escape the scan. Sol and Gemini's selected assistant-content fields are empty throughout, while Grok's permit the named-call comparisons. Private notes and returned explanations also differ in availability and fidelity. We therefore report source-qualified propositions, not cross-model deception or monitorability scores. A monitor evaluation needs a declared observation channel, prediction target and held-out criterion \citep{korbak2025monitorability}.
